\documentclass[11pt,a4paper]{article}
\usepackage[T1]{fontenc}
\usepackage[utf8]{inputenc}
\usepackage[french]{babel}
\usepackage{geometry}
\usepackage{amsmath,amssymb}
\usepackage{enumitem}
\usepackage{hyperref}
\geometry{margin=2.5cm}

\title{Trois exercices sur l'allocation optimale en sondage stratifié}
\author{}
\date{}

\begin{document}
\maketitle

\section*{Consigne générale}
Pour chaque exercice, on suppose que l'on cherche une allocation optimale dans un sondage stratifié.  
Les données numériques sont fournies dans des fichiers Excel séparés afin de pouvoir être exploitées directement dans une application de calcul.

\bigskip
\noindent Fichiers de données associés :
\begin{itemize}[leftmargin=2em]
    \item \texttt{exercice\_1\_sondage\_stratifie.xlsx}
    \item \texttt{exercice\_2\_sondage\_stratifie.xlsx}
    \item \texttt{exercice\_3\_sondage\_stratifie.xlsx}
\end{itemize}

\section{Exercice 1 : allocation optimale pour une seule variable}
On considère une population divisée en quatre strates.

Le fichier \texttt{exercice\_1\_sondage\_stratifie.xlsx} contient, pour chaque strate :
\begin{itemize}[leftmargin=2em]
    \item l'effectif \texttt{Nh},
    \item les quasi-variances \texttt{S2\_T}, \texttt{S2\_Y} et \texttt{S2\_Z}.
\end{itemize}

On veut déterminer les entiers positifs $n_h$ qui minimisent
\[
n = \sum_h n_h
\]
sous la contrainte principale
\[
\mathrm{Var}(T) \leq \alpha_T.
\]
On vérifiera ensuite que les allocations obtenues respectent aussi les bornes
\[
\mathrm{Var}(Y) \leq \alpha_Y
\qquad \text{et} \qquad
\mathrm{Var}(Z) \leq \alpha_Z.
\]

\textbf{Travail demandé :}
\begin{enumerate}[leftmargin=2em]
    \item Calculer l'allocation optimale.
    \item Donner les fractions de sondage $f_h$.
    \item Calculer les trois variances finales.
    \item Interpréter la strate la plus coûteuse en taille d'échantillon.
\end{enumerate}

\section{Exercice 2 : allocation simultanée pour trois variables}
Le fichier \texttt{exercice\_2\_sondage\_stratifie.xlsx} contient cinq strates et les informations nécessaires à l'optimisation.

On cherche maintenant à résoudre :
\[
\min \sum_h n_h
\]
sous les contraintes
\[
\mathrm{Var}(T) \leq \alpha_T, \qquad
\mathrm{Var}(Y) \leq \alpha_Y, \qquad
\mathrm{Var}(Z) \leq \alpha_Z.
\]

\textbf{Travail demandé :}
\begin{enumerate}[leftmargin=2em]
    \item Déterminer les $n_h$ optimaux entiers positifs.
    \item Calculer la taille totale optimale $n$.
    \item Comparer les contributions à la variance par strate.
    \item Commenter l'influence des strates les plus dispersées.
\end{enumerate}

\section{Exercice 3 : comparaison entre allocation proportionnelle et allocation optimale}
Le fichier \texttt{exercice\_3\_sondage\_stratifie.xlsx} contient une population plus contrastée.

On demande de comparer :
\begin{itemize}[leftmargin=2em]
    \item une allocation proportionnelle,
    \item une allocation optimale sous contraintes de variance.
\end{itemize}

\textbf{Travail demandé :}
\begin{enumerate}[leftmargin=2em]
    \item Calculer une allocation proportionnelle de référence.
    \item Calculer l'allocation optimale.
    \item Comparer les tailles totales obtenues.
    \item Comparer les variances finales et mesurer le gain apporté par l'optimisation.
\end{enumerate}

\section*{Rappel de la formule de variance}
Pour une variable quelconque $X$, la variance du sondage stratifié est donnée par
\[
\mathrm{Var}(X)
=
\sum_h
\left(\frac{N_h}{N}\right)^2
\left(\frac{1-f_h}{n_h}\right)
S_h^2(X),
\qquad
f_h = \frac{n_h}{N_h}.
\]

\end{document}
